\section{Appendix A. Mathematical Derivations}

\subsection{A.1 Kernel Attention Derivation}

Starting from softmax attention:

\[
A_i = \frac{\sum_j \kappa(q_i, k_j) v_j}{\sum_j \kappa(q_i, k_j)}
\tag{A.1}
\]

If $\kappa(q, k) = \phi(q)^\top \phi(k)$:

\[
A_i = \frac{\phi(q_i)^\top \sum_j \phi(k_j) v_j^\top}{\phi(q_i)^\top \sum_j \phi(k_j)}
\tag{A.2}
\]

The complexity is $\mathcal{O}(L F d_h)$.

\subsection{A.2 Bochner's Theorem}

For a shift-invariant kernel $\kappa(x - y)$:

\[
\kappa(x - y) = \int p(\omega) e^{i \omega^\top (x - y)} d\omega
\tag{A.3}
\]

For the Gaussian kernel:

\[
\kappa(\Delta) = \exp\!\left(-\frac{\|\Delta\|^2}{2}\right) \implies p(\omega) = \mathcal{N}(0, I)
\tag{A.4}
\]

\subsection{A.3 RFF Estimator}

\[
\kappa(x - y) = \mathbb{E}_\omega[\cos(\omega^\top (x - y))] = \mathbb{E}_\omega[\cos(\omega^\top x) \cos(\omega^\top y) + \sin(\omega^\top x) \sin(\omega^\top y)]
\tag{A.5}
\]

Sampling $F$ frequencies $\omega_1, \dots, \omega_F$:

\[
\hat{\kappa}(x - y) = \frac{1}{F} \sum_i [\cos(\omega_i^\top x) \cos(\omega_i^\top y) + \sin(\omega_i^\top x) \sin(\omega_i^\top y)]
\tag{A.6}
\]

Define:

\[
\phi(x) = \frac{1}{\sqrt{F}} [\cos(\omega_1^\top x), \dots, \cos(\omega_F^\top x), \sin(\omega_1^\top x), \dots, \sin(\omega_F^\top x)]
\tag{A.7}
\]

Then $\hat{\kappa}(x - y) = \phi(x)^\top \phi(y)$.

\subsection{A.4 Variance of the Estimator}

For iid sampling:

\[
\text{Var}[\hat{\kappa}] = \frac{1}{F} \text{Var}[\cos(\omega^\top \Delta)] = \mathcal{O}(1/F)
\tag{A.8}
\]

For orthogonalized sampling:

\[
\text{Var}[\hat{\kappa}] = \mathcal{O}(1/F^2)
\tag{A.9}
\]

\subsection{A.5 Causal Cumsum}

For decoders, we need cumulative sums:

\[
\text{KV}_i = \sum_{j \leq i} \phi(k_j) v_j^\top
\tag{A.10}
\]

This is computed in $\mathcal{O}(L)$ via \texttt{torch.cumsum} on the outer product.

\textbf{Memory.} The outer product has shape $(B, H, L, 2F, d_h)$. For 
$L = 128$, this is ~1.6$\times$10  elements at FP32 = 640 MB.

\subsection{A.6 Preservation Gate}

\[
\alpha = \sigma(w), \quad \text{out} = \alpha A + (1 - \alpha) x
\tag{A.11}
\]

Gradient w.r.t. $w$:

\[
\frac{\partial \text{out}}{\partial w} = \alpha(1 - \alpha)(A - x)
\tag{A.12}
\]

The gate learns to increase $\alpha$ when $A$ has larger magnitude than $x$, and 
decrease when $x$ is more important.

\subsection{A.7 Parameter Count}

\begin{table}[H]
\centering
\small
\begin{tabular}{ll}
\toprule
\textbf{Component} & \textbf{Params} \\
\midrule
Q/K/V/O (frozen) & $4 d^2$ \\
LoRA ($\times$4) & $4 \cdot 2 r_L d$ \\
Random projection (per head) & $H d_h F$ \\
LayerNorm & $2 d_h$ \\
Gates & 2 \\
\bottomrule
\end{tabular}
\end{table}

For $d = 768$, $H = 12$, $d_h = 64$, $F = 256$, $r_L = 16$:

\begin{table}[H]
\centering
\small
\begin{tabular}{ll}
\toprule
\textbf{Component} & \textbf{Params} \\
\midrule
Q/K/V/O & 2,359,296 \\
LoRA & 98,304 \\
Projection & 196,608 \\
LN & 128 \\
Gates & 2 \\
\textbf{Total} & 2,654,338 \\
\textbf{vs baseline} & 2,359,296 \\
\textbf{Increase} & 12.5\% \\
\bottomrule
\end{tabular}
\end{table}

